% Source context: arXiv:2203.12563v3, REsubmission.tex, line 1584;
% arXiv:1010.3732, Section II.F.2. The quantitative bounds below are
% auxiliary physical-enlargement results, not the full classification theorem.
% The linked declarations and axiom guards were compiler-checked in the
% arbitrary-physical mixed-path package.

\section{Physical isometries and periodic gaps}
\label{sec:mpo_isometric_gap_transport}

The physical-support reduction in
\cite[Section~5]{GarreRubio2023MPOSym} restricts each endpoint to the
subspace spanned by its tensor. To compare Hamiltonians before and after
an isometric inclusion, the unused physical space must have positive
energy. Let $E:\C^d\to\C^m$ satisfy $E^\dagger E=I$, put
$Q=EE^\dagger$, $T=E\otimes E$, and $W_N=E^{\otimes N}$.
For a positive two-site interaction $h$, use the extension of
Definition~\ref{def:spt_isometric_interaction_extension}:
\begin{align}
    h' & =ThT^\dagger+B, & B & =I-Q\otimes Q.
    \label{eq:mpo_isometry_extension}
\end{align}
For $N\geq2$, write $h_i$, $h'_i$, and $B_i$ for the respective
translates on sites $i,i+1$ modulo $N$, and put
$H_N=\sum_{i=0}^{N-1}h_i$, $H'_N=\sum_{i=0}^{N-1}h'_i$, and
$P_N=\sum_{i=0}^{N-1}B_i$. Neither the $h_i$ nor the $h'_i$ are
assumed to commute.

The conjugated term $C=ThT^\dagger$ in
(\ref{eq:mpo_isometry_extension}) has components
\begin{align}
    C_{\alpha_0\alpha_1,\beta_0\beta_1}
    =\sum_{a_0,a_1,b_0,b_1=0}^{d-1}
    E_{\alpha_0a_0}E_{\alpha_1a_1}
    h_{a_0a_1,b_0b_1}
    \overline{E_{\beta_0b_0}E_{\beta_1b_1}}.
    \label{eq:mpo_isometry_components}
\end{align}
The north legs are output indices $\alpha_0,\alpha_1$ and the south
legs are input indices $\beta_0,\beta_1$, each of dimension $m$.
The four internal legs have dimension $d$; the lower tensors are
$E^\dagger$, so their coefficients are complex conjugated as in
(\ref{eq:mpo_isometry_components}).
% TENKZ-ISOMETRIC-GAP-BEGIN
% Formula: C = (E tensor E) h (E tensor E)^dagger.
% Ink: upper boxes E_{alpha_i,a_i}; center h_{a_0 a_1,b_0 b_1};
% lower boxes (E^dagger)_{b_i,beta_i}=conjugate(E_{beta_i,b_i}).
% Contracted indices a_0,a_1,b_0,b_1 each range over Fin d.
% Open alpha_0,alpha_1 (north outputs) and beta_0,beta_1 (south inputs)
% each range over Fin m. No virtual bonds and no inter-isometry edge.
% Boundary signature on each panel: phys:n, phys:n, phys:s, phys:s.
\begin{tenkzequation}
    \begin{tenkz}[rows={wire}, cols=2, bonds=none]
        \tn[at=(1,1), name=C, skin=box, wide=2, size=l,
            ports={90@1:physical:$\alpha_0$, 90@2:physical:$\alpha_1$,
                    270@1:physical:$\beta_0$, 270@2:physical:$\beta_1$}]{C}
    \end{tenkz}
    $=$
    \begin{tenkz}[rows={wire,wire,wire}, cols=2, bonds=none]
        \tn[at=(1,1), name=Ezero, skin=box,
            ports={90:physical:$\alpha_0$, 270:physical}]{E}
        \tn[at=(1,2), name=Eone, skin=box,
            ports={90:physical:$\alpha_1$, 270:physical}]{E}
        \tn[at=(2,1), name=h, skin=box, wide=2, size=l,
            ports={90@1:physical, 90@2:physical,
                    270@1:physical, 270@2:physical}]{h}
        \tn[at=(3,1), name=Fzero, skin=box,
            ports={90:physical, 270:physical:$\beta_0$}]{E^\dagger}
        \tn[at=(3,2), name=Fone, skin=box,
            ports={90:physical, 270:physical:$\beta_1$}]{E^\dagger}
        \tnwire{Ezero.270}{h.90@1}
        \tnwire{Eone.270}{h.90@2}
        \tnwire{h.270@1}{Fzero.90}
        \tnwire{h.270@2}{Fone.90}
    \end{tenkz}
\end{tenkzequation}
% TENKZ-ISOMETRIC-GAP-END

\subsection{Commuting support penalties}

\begin{theorem}[Commutation of tensor-square translates]
    \label{thm:mpo_isometry_support_commute}
    \lean{MPOTensor.embed_twoSiteSectorProjection_commute,
        MPSTensor.isometricInteractionExtension_zero_translate_commute}
    \leanok
    For any one-site matrix $R$ and $N\geq2$, all periodic translates
    of $R\otimes R$ commute. Thus the translates of $I-Q\otimes Q$
    commute. The commutation assertion requires neither $R$ to be a
    projection nor $E$ to be an isometry when $Q=EE^\dagger$.
\end{theorem}
\begin{proof}\leanok
    A translate is the tensor product of $R$ on its two cyclic sites
    and $I$ elsewhere. At each site, the factors from any two translates
    belong to $\{R,I\}$ and commute. Multiplication of tensor products
    gives commutation on the chain. Taking complements preserves it.
    The same argument includes $N=2$, when both translates cover the
    whole ring.
\end{proof}

\begin{theorem}[Energy outside the included chain]
    \label{thm:mpo_isometry_inactive_bound}
    \lean{MPSTensor.isometricInteractionExtension_zero_ker_eq_range,
        MPSTensor.isometricInteractionExtension_posSemidef,
        MPSTensor.isometricInteractionExtension_inactive_energy_lower_bound}
    \leanok
    \uses{def:spt_isometric_interaction_extension}
    If $E^\dagger E=I$, $h\geq0$, and $N\geq2$, then $h'\geq0$ and
    \begin{align}
        \ker P_N & =\operatorname{ran}W_N,
        \label{eq:mpo_isometry_penalty_kernel}                 \\
        \|y\|^2  & \leq\operatorname{Re}\langle H'_Ny,y\rangle
        \quad (y\perp\operatorname{ran}W_N).
        \label{eq:mpo_isometry_inactive_energy}
    \end{align}
\end{theorem}
\begin{proof}\leanok
    \uses{thm:mpo_isometry_support_commute,
        thm:spt_isometric_interaction_ground_space,
        thm:spt_sector_bond_gap,thm:spt_isometric_parent_interaction,
        thm:spt_periodic_hamiltonian_order}
    The interaction $B$ is the isometric extension of zero and is an
    orthogonal projection. Apply the exact kernel transport for
    commuting projections to the zero interaction and $B$; the original
    zero Hamiltonian has the whole chain as its kernel. This gives
    (\ref{eq:mpo_isometry_penalty_kernel}). The commuting-projection
    spectral bound gives $\operatorname{spec}P_N\subseteq
        \{0\}\cup[1,\infty)$, hence
    $\langle P_Ny,y\rangle\geq\|y\|^2$ on its kernel complement.
    Positivity of $h$ gives $0\leq B\leq h'$ and $P_N\leq H'_N$.
    Comparing quadratic forms proves
    (\ref{eq:mpo_isometry_inactive_energy}). Only the auxiliary
    penalties $B_i$ have been required to commute.
\end{proof}

\subsection{Orthogonal decomposition along an isometry}

Let $U:\mathcal E\to\mathcal F$ be an isometry between finite-dimensional
complex inner-product spaces. Suppose $K$ and $L$ are linear operators
on $\mathcal E$ and $\mathcal F$, respectively, with $LU=UK$.
Write $q_L(v)=\operatorname{Re}\langle Lv,v\rangle$ and
$q_K(x)=\operatorname{Re}\langle Kx,x\rangle$.
Every $v\in\mathcal F$ has the orthogonal decomposition
\begin{align}
    x       & =U^\dagger v,     & y & =v-Ux, & U^\dagger y & =0,
    \label{eq:mpo_isometry_decomposition}                        \\
    \|v\|^2 & =\|x\|^2+\|y\|^2.
    \label{eq:mpo_isometry_norm_split}
\end{align}

\begin{theorem}[Kernel transport from a strict complementary bound]
    \label{thm:mpo_isometry_abstract_kernel}
    \lean{LinearIsometry.ker_eq_map_of_intertwines_of_inactive_bound}
    \leanok
    If $c>0$ and $c\|y\|^2\leq q_L(y)$ whenever $U^\dagger y=0$,
    then $\ker L=U(\ker K)$. No symmetry or positivity assumption on
    $L$ or $K$ is needed.
\end{theorem}
\begin{proof}\leanok
    If $Lv=0$, use (\ref{eq:mpo_isometry_decomposition}). Then
    $Ly=-UKx$, so orthogonality gives $q_L(y)=0$. The strict lower
    bound forces $y=0$. Thus $v=Ux$, and $UKx=Lv=0$ implies $Kx=0$
    by injectivity of $U$. Conversely, $Kx=0$ gives $LUx=UKx=0$.
\end{proof}

\begin{theorem}[Combining quadratic-form bounds]
    \label{thm:mpo_isometry_abstract_quadratic}
    \lean{LinearIsometry.re_inner_ge_of_intertwines}
    \leanok
    Suppose $L$ is self-adjoint, $\delta,c\in\R$, and
    \begin{align}
        \delta\|x\|^2 & \leq q_K(x)
                      &             & (x\perp\ker K),\notag  \\
        c\|y\|^2      & \leq q_L(y)
                      &             & (U^\dagger y=0).\notag
    \end{align}
    Then $\min\{\delta,c\}\|v\|^2\leq q_L(v)$ for
    $v\perp\ker L$. Neither constant is required to be nonnegative.
\end{theorem}
\begin{proof}\leanok
    For $v\perp\ker L$, the vector $x=U^\dagger v$ lies in
    $(\ker K)^\perp$, because $U(\ker K)\subseteq\ker L$.
    Intertwining gives $\langle LUx,y\rangle=0$ in
    (\ref{eq:mpo_isometry_decomposition}). Self-adjointness gives
    $\langle Ly,Ux\rangle=0$ as well. Therefore
    \begin{align}
        q_L(v) & =q_K(x)+q_L(y)
        \geq\delta\|x\|^2+c\|y\|^2\notag                         \\
               & \geq\min\{\delta,c\}\bigl(\|x\|^2+\|y\|^2\bigr)
        =\min\{\delta,c\}\|v\|^2.\notag
    \end{align}
    The last equality is (\ref{eq:mpo_isometry_norm_split}).
\end{proof}

\begin{theorem}[Combining a norm gap with a complementary energy bound]
    \label{thm:mpo_isometry_abstract_norm}
    \lean{LinearIsometry.norm_gap_of_intertwines}
    \leanok
    Suppose $K,L\geq0$, $\delta\geq0$, and
    $\delta\|x\|\leq\|Kx\|$ for $x\perp\ker K$.
    If $c\|y\|^2\leq q_L(y)$ whenever $U^\dagger y=0$, then
    $\min\{\delta,c\}\|v\|\leq\|Lv\|$ for $v\perp\ker L$.
    No sign assumption on $c$ is required.
\end{theorem}
\begin{proof}\leanok
    \uses{thm:mpo_isometry_abstract_quadratic}
    For a positive operator, the norm-gap hypothesis with
    $\delta\geq0$ gives $q_K(x)\geq\delta\|x\|^2$ on the kernel
    complement. Apply Theorem~\ref{thm:mpo_isometry_abstract_quadratic}
    and the bound $q_L(v)\leq\|Lv\|\|v\|$.
    When $v\neq0$, division by $\|v\|$ gives the assertion;
    for $v=0$ it is immediate.
\end{proof}

\subsection{Periodic ground spaces and gap bounds}

\begin{theorem}[Exact ground-space transport for positive interactions]
    \label{thm:mpo_isometry_positive_kernel}
    \lean{MPSTensor.isometricInteractionExtension_ker_eq_map_of_posSemidef,
        MPSTensor.isometricInteractionExtension_ker_eq_span_of_posSemidef}
    \leanok
    \uses{def:spt_isometric_interaction_extension}
    For $E^\dagger E=I$, $h\geq0$, and $N\geq2$,
    \begin{align}
        \ker H'_N=W_N(\ker H_N).
        \label{eq:mpo_isometry_kernel_transport}
    \end{align}
    In particular, if $\ker H_N=\spn_{\C}\{\psi\}$, then
    $\ker H'_N=\spn_{\C}\{W_N\psi\}$. This holds even for
    $\psi=0$; a nonzero ground vector is not a hypothesis.
\end{theorem}
\begin{proof}\leanok
    \uses{thm:spt_isometric_interaction_zero_modes,
        thm:mpo_isometry_inactive_bound,thm:mpo_isometry_abstract_kernel}
    The tensor power $W_N$ is an isometry, and the periodic
    intertwining identity is $H'_NW_N=W_NH_N$.
    Theorem~\ref{thm:mpo_isometry_inactive_bound} supplies the strict
    complementary bound with $c=1$. Apply
    Theorem~\ref{thm:mpo_isometry_abstract_kernel}. The image of a
    span is the span of the image, giving the final assertion.
\end{proof}

\begin{theorem}[Gap transport with a unit unused-space penalty]
    \label{thm:mpo_isometry_positive_gap}
    \lean{MPSTensor.isometricInteractionExtension_quadratic_gap,
        MPSTensor.isometricInteractionExtension_norm_gap,
        MPSTensor.isometricInteractionExtension_spectrum_gap}
    \leanok
    \uses{def:spt_isometric_interaction_extension}
    Suppose $E^\dagger E=I$, $h\geq0$, and $N\geq2$.
    A quadratic-form bound $\delta\|x\|^2\leq
        \operatorname{Re}\langle H_Nx,x\rangle$ on $(\ker H_N)^\perp$
    gives
    \begin{align}
        \min\{\delta,1\}\|v\|^2
        \leq\operatorname{Re}\langle H'_Nv,v\rangle
        \quad (v\perp\ker H'_N).
        \label{eq:mpo_isometry_quadratic_gap}
    \end{align}
    If $\delta\geq0$, a norm gap $\delta\|x\|\leq\|H_Nx\|$
    on $(\ker H_N)^\perp$ gives
    $\min\{\delta,1\}\|v\|\leq\|H'_Nv\|$ on $(\ker H'_N)^\perp$.
    Finally, spectral separation
    $\operatorname{spec}H_N\subseteq\{0\}\cup[\delta,\infty)$ gives
    \begin{align}
        \operatorname{spec}H'_N
        \subseteq[0,\infty)\cap
        \bigl(\{0\}\cup[\min\{\delta,1\},\infty)\bigr).
        \label{eq:mpo_isometry_spectral_gap}
    \end{align}
    The quadratic-form and spectral assertions allow arbitrary real
    $\delta$. None of these assertions requires $0$ to be an eigenvalue.
\end{theorem}
\begin{proof}\leanok
    \uses{thm:spt_isometric_interaction_zero_modes,
        thm:spt_periodic_hamiltonian_order,thm:mpo_isometry_inactive_bound,
        thm:mpo_isometry_abstract_quadratic,thm:mpo_isometry_abstract_norm}
    Positivity of $h$ and $h'$ makes both periodic sums positive.
    Combine $H'_NW_N=W_NH_N$ with
    (\ref{eq:mpo_isometry_inactive_energy}) and apply the abstract
    quadratic-form or norm theorem with $U=W_N$ and $c=1$.
    For the spectral assertion, finite-dimensional spectral separation
    first gives the quadratic-form bound on $(\ker H_N)^\perp$.
    Transfer it by (\ref{eq:mpo_isometry_quadratic_gap}) and apply
    the converse spectral implication to the positive operator $H'_N$.
\end{proof}

In particular, a positive bound $\delta$ uniform in $N\geq2$ remains
uniform after physical inclusion, with constant $\min\{\delta,1\}$.
For canonical two-site parent interactions,
Theorem~\ref{thm:spt_isometric_parent_interaction} identifies $h'$ with
the parent interaction of the physically included tensor. Thus the gap
estimate applies when restoring a larger physical alphabet after a
support reduction. Constructing the endpoint support coordinates,
proving covariance under a common MPO symmetry along an interpolation,
and classifying phases with degenerate ground spaces remain separate
mathematical assertions.

% Source context: arXiv:2203.12563v3, REsubmission.tex, lines 1584--1601
% and 1687--1692. The arbitrary-alphabet continuation is compiler-checked.
\section{Mixed paths in arbitrary physical alphabets}
\label{sec:mpo_arbitrary_physical_path}

Let $A_p=(A_p^i)_{i=0}^{d_p-1}$ have bond dimension $D_p$, for
$p=0,1$, and put $D=D_0+D_1$. The physical dimensions $d_p$ need
not equal $D_p^2$. Read $A_p$ as a matrix
$\mathscr A_p:\C^{D_p^2}\to\C^{d_p}$ and take its polar
factorization $\mathscr A_p=F_pP_p$. Let $B_p$ be the tensor formed
from the rows of $P_p$, with physical dimension $D_p^2$ and bond
dimension $D_p$, as in Definition~\ref{def:ldpolar_tensor}.
Thus $A_p=F_pB_p$, where multiplication denotes physical mixing.
This realizes the physical-support reduction preceding the mixed
interpolation in \cite[Section~5]{GarreRubio2023MPOSym}.

\subsection{Polar support and full endpoint embeddings}

\begin{theorem}[Injectivity of the positive polar tensor]
    \label{thm:mpo_polar_support_injective}
    \lean{MPSTensor.isInjective_polarPosTensor}
    \leanok
    \uses{def:ldpolar_tensor,def:injective}
    If $A_p$ is one-site injective, then $B_p$ is one-site injective.
\end{theorem}
\begin{proof}\leanok
    \uses{thm:ldpolar_tensor_decomposition}
    Every letter $A_p^i=\sum_k(F_p)_{ik}B_p^k$ belongs to the span
    of the letters of $B_p$. Since the letters of $A_p$ span
    $\MN{D_p}$, the letters of $B_p$ span that algebra as well.
\end{proof}

Let $I_p:\C^{D_p^2}\to\C^{D^2}$ include the $p$th diagonal
physical sector, and put
$\mathcal H=\C^{D^2}\oplus\C^{d_0}\oplus\C^{d_1}$ and
$Jw=(w,0,0)$. The full original physical spaces enter $\mathcal H$
through the maps of Definition~\ref{def:spt_common_physical_endpoints}:
\begin{align}
    E_0v & =(I_0F_0^\dagger v,(I-F_0F_0^\dagger)v,0),\notag \\
    E_1v & =(I_1F_1^\dagger v,0,(I-F_1F_1^\dagger)v).
    \label{eq:mpo_arbitrary_full_embeddings}
\end{align}
When $A_p$ is injective, Theorems~\ref{thm:ldpolar_tensor_injective}
and~\ref{thm:spt_common_physical_endpoints} give
$F_p^\dagger F_p=I$ and $E_p^\dagger E_p=I$.
The original unused directions are retained in the separate
$\C^{d_p}$ summands. The disjoint coordinate sectors also give
$E_0^\dagger E_1=0$, by
Theorem~\ref{thm:spt_common_physical_endpoint_orthogonality}.
These assertions require no symmetry of the endpoint tensors.

\begin{definition}[The mixed tensor in the common physical space]
    \label{def:mpo_arbitrary_physical_tensor}
    \lean{MPSTensor.MPOSymmetry.arbitraryPhysicalMixedInterpolation}
    \leanok
    \uses{def:mpo_mixed_endpoint_tensor,def:ldpolar_tensor,
        def:spt_common_physical_endpoints}
    Let $\mathcal B(\gamma)$ be the mixed interpolation of $B_0,B_1$
    from Definition~\ref{def:mpo_mixed_endpoint_tensor}. Define
    \begin{align}
        \widetilde A(\gamma)=J\mathcal B(\gamma).
        \label{eq:mpo_arbitrary_tensor}
    \end{align}
    Its physical dimension is $D^2+d_0+d_1$ and its bond dimension is $D$.
\end{definition}

\begin{theorem}[Continuity and interior injectivity of the enlarged tensor]
    \label{thm:mpo_arbitrary_tensor_continuity}
    \lean{MPSTensor.MPOSymmetry.continuous_arbitraryPhysicalMixedInterpolation,
        MPSTensor.MPOSymmetry.isInjective_arbitraryPhysicalMixedInterpolation}
    \leanok
    \uses{def:mpo_arbitrary_physical_tensor,def:injective}
    For arbitrary $A_0,A_1$, the tensor $\widetilde A(\gamma)$ is
    continuous for $\gamma\in\R$. If both tensors are one-site
    injective, then $\widetilde A(\gamma)$ is injective for
    $0<\gamma<1$.
\end{theorem}
\begin{proof}\leanok
    \uses{thm:mpo_mixed_interior,thm:mpo_polar_support_injective,
        thm:spt_common_physical_endpoints}
    The mixed tensor depends continuously on $\gamma$, and physical
    mixing by the fixed matrix $J$ is continuous. Injectivity of
    $B_0,B_1$ gives interior injectivity of $\mathcal B(\gamma)$.
    Since $J^\dagger J=I$, applying $J^\dagger$ to the physical
    letters recovers $\mathcal B(\gamma)$ and preserves their full span.
\end{proof}

\begin{theorem}[Original endpoint vectors in orthogonal physical spaces]
    \label{thm:mpo_arbitrary_endpoint_vectors}
    \lean{MPSTensor.MPOSymmetry.mixedEndpointLeftTensor_eq_rotatePhysical,
        MPSTensor.MPOSymmetry.mixedEndpointRightTensor_eq_rotatePhysical,
        MPSTensor.MPOSymmetry.rotatePhysical_mixedEndpointLeftTensor_polarPosTensor,
        MPSTensor.MPOSymmetry.rotatePhysical_mixedEndpointRightTensor_polarPosTensor,
        MPSTensor.MPOSymmetry.mpv_arbitraryPhysicalMixedInterpolation_zero,
        MPSTensor.MPOSymmetry.mpv_arbitraryPhysicalMixedInterpolation_one}
    \leanok
    \uses{def:mpo_arbitrary_physical_tensor,def:mpo_mixed_first_endpoint,
        def:mpo_mixed_endpoint_swap}
    Extending a square-alphabet tensor by zero into its diagonal
    physical sector is exactly physical mixing by $I_p$.
    In particular, the auxiliary small-bond tensors
    $\widehat B_p=I_pB_p$ satisfy $J\widehat B_p=E_pA_p$ whenever
    $A_p$ is injective. For every $N>0$, injectivity of the corresponding
    endpoint gives
    \begin{align}
        \ket{V^{(N)}(\widetilde A(0))}
         & =\ket{V^{(N)}(E_0A_0)},\notag \\
        \ket{V^{(N)}(\widetilde A(1))}
         & =\ket{V^{(N)}(E_1A_1)}.
        \label{eq:mpo_arbitrary_endpoint_vectors}
    \end{align}
    The equality at $p$ requires injectivity only of $A_p$.
\end{theorem}
\begin{proof}\leanok
    \uses{thm:spt_common_physical_polar_endpoints,
        thm:ldpolar_tensor_decomposition,
        thm:mpo_mixed_periodic_line_continuity}
    The coordinate inclusion agrees with the original tensor on its
    sector and vanishes elsewhere. The identities
    $E_pF_p=JI_p$ and $A_p=F_pB_p$ give
    $E_pA_p=JI_pB_p=J\widehat B_p$.
    The mixed endpoint tensors have block forms
    \begin{align}
        \mathcal B(0)^\mu
         & =\begin{pmatrix}\widehat B_0^\mu&0\\ R^\mu&0\end{pmatrix},
         & \mathcal B(1)^\mu
         & =\begin{pmatrix}0&S^\mu\\0&\widehat B_1^\mu\end{pmatrix}.
        \label{eq:mpo_arbitrary_endpoint_blocks}
    \end{align}
    The off-diagonal blocks do not contribute to the trace of any
    nonempty word. Hence their periodic vectors equal those of
    $\widehat B_0$ and $\widehat B_1$, respectively. Physical mixing
    by $J$ acts on every physical index of these vector equalities and
    gives (\ref{eq:mpo_arbitrary_endpoint_vectors}).
\end{proof}

The off-diagonal blocks in (\ref{eq:mpo_arbitrary_endpoint_blocks})
are retained by the actual tensor path. The endpoint identities
(\ref{eq:mpo_arbitrary_endpoint_vectors}) compare periodic vectors;
they do not identify the raw bond-$D$ tensors with the bond-$D_p$
original tensors.

\subsection{The continuous positive interaction}

\begin{theorem}[Coordinate representation of a periodic interaction]
    \label{thm:mpo_periodic_interaction_coordinates}
    \lean{MPSTensor.periodicLocalInteractionES_eq_toEuclideanLin_embedLocalOperator,
        MPSTensor.periodicInteractionHamiltonianES_eq_toEuclideanLin_sum_embedLocalOperator,
        MPSTensor.periodicInteractionHamiltonianES_eq_toEuclideanLin_interactionHamiltonian}
    \leanok
    For any local matrix $a$ on $(\C^d)^{\otimes R}$, with $R\leq N$,
    its action on a cyclic window agrees with the operator defined by applying $a$
    to that window and fixing the exterior configuration. Their periodic
    sums therefore agree. In particular, for $R=2$ and $N\geq2$,
    the sum of the Euclidean local operators is the operator with
    configuration-basis matrix $\sum_{i=0}^{N-1}a_i$.
    No positivity or projection assumption on $a$ is needed.
\end{theorem}
\begin{proof}\leanok
    For a configuration $\sigma$, let $\sigma|_i$ be its ordered
    cyclic window and let $\sigma[i\leftarrow\tau]$ replace that
    window by $\tau$. Both local operators have the coefficient formula
    \begin{align}
        (a_iv)_\sigma
        =\sum_{\tau\in\{0,\ldots,d-1\}^{R}}
        a_{\sigma|_i,\tau}\,v_{\sigma[i\leftarrow\tau]}.
        \label{eq:mpo_periodic_matrix_action}
    \end{align}
    The cyclic window and its complement partition the sites, so their
    coordinate permutation extracts and replaces precisely these indices.
    Sum (\ref{eq:mpo_periodic_matrix_action}) over the cyclic origins.
\end{proof}

\begin{definition}[Enlarged mixed local interaction]
    \label{def:mpo_arbitrary_physical_interaction}
    \lean{MPSTensor.MPOSymmetry.mixedEndpointInteractionMatrix,
        MPSTensor.MPOSymmetry.toEuclideanLin_mixedEndpointInteractionMatrix,
        MPSTensor.MPOSymmetry.arbitraryPhysicalMixedInteraction}
    \leanok
    \uses{def:mpo_mixed_extended_support,def:mpo_arbitrary_physical_tensor,
        def:spt_isometric_interaction_extension}
    For any square-alphabet tensors $C_0,C_1$, let
    $h_{C_0,C_1}(\gamma)$ be the configuration-basis matrix of
    $I-\Pi_{\mathcal S'_\gamma}$ for their extended mixed two-site
    support. It represents exactly that Euclidean operator. Set
    $h_\gamma=h_{B_0,B_1}(\gamma)$, $T=J\otimes J$, and define
    \begin{align}
        \widetilde h_\gamma
                     & =Th_\gamma T^\dagger+I-TT^\dagger,\notag                                             \\
        H_{\gamma,N} & =\sum_{i=0}^{N-1}(h_\gamma)_i,
                     & \widetilde H_{\gamma,N}                  & =\sum_{i=0}^{N-1}(\widetilde h_\gamma)_i.
        \label{eq:mpo_arbitrary_interaction}
    \end{align}
    The periodic sums are defined for $N\geq2$.
\end{definition}

\begin{theorem}[Positivity and continuity through the endpoints]
    \label{thm:mpo_arbitrary_interaction_continuity}
    \lean{MPSTensor.MPOSymmetry.mixedEndpointInteractionMatrix_posSemidef,
        MPSTensor.MPOSymmetry.continuous_mixedEndpointInteractionMatrix,
        MPSTensor.MPOSymmetry.arbitraryPhysicalMixedInteraction_posSemidef,
        MPSTensor.MPOSymmetry.continuous_arbitraryPhysicalMixedInteraction}
    \leanok
    \uses{def:mpo_arbitrary_physical_interaction,def:injective}
    For arbitrary $C_0,C_1$ and real $\gamma$,
    $h_{C_0,C_1}(\gamma)\geq0$; for arbitrary $A_0,A_1$,
    $\widetilde h_\gamma\geq0$.
    If $D_0,D_1>0$ and $C_0,C_1$ are injective, then
    $h_{C_0,C_1}(\gamma)$ is continuous on $\R$.
    Thus, if $D_0,D_1>0$ and $A_0,A_1$ are one-site injective,
    $\widetilde h_\gamma$ is continuous on $\R$, including at $0$ and $1$.
\end{theorem}
\begin{proof}\leanok
    \uses{thm:mpo_mixed_endpoint_projector_comparison,
        thm:mpo_mixed_support_continuity,thm:mpo_polar_support_injective,
        thm:mpo_isometry_inactive_bound}
    The complement of an orthogonal projection is positive.
    Passing to matrix coordinates preserves positivity and continuity
    in finite dimension. Under the stated injectivity hypotheses,
    the positive polar tensors are injective, so the continuous-support
    theorem applies to $h_\gamma$.
    Since $T$ is fixed, (\ref{eq:mpo_arbitrary_interaction}) is an
    affine continuous function of $h_\gamma$. Positivity of the
    isometric extension gives $\widetilde h_\gamma\geq0$.
\end{proof}

\begin{theorem}[Canonical parent equality in the open interval]
    \label{thm:mpo_arbitrary_interior_parent}
    \lean{MPSTensor.MPOSymmetry.arbitraryPhysicalMixedInteraction_eq_parent_of_mem_Ioo}
    \leanok
    \uses{def:mpo_arbitrary_physical_interaction,def:parent_interaction}
    For arbitrary original tensors $A_0,A_1$ and $0<\gamma<1$,
    $\widetilde h_\gamma$ equals the canonical two-site parent
    interaction of $\widetilde A(\gamma)$.
\end{theorem}
\begin{proof}\leanok
    \uses{thm:mpo_mixed_support_continuity,thm:spt_isometric_parent_interaction}
    In the open interval the insertion in the mixed tensor is invertible,
    so its extended two-site support is the canonical one. The
    canonical parent of a tensor physically included by $J$ is its
    isometric interaction extension. Apply these two identities.
\end{proof}

\subsection{Ground lines and a volume-uniform gap}

\begin{theorem}[The exact nonzero periodic ground line]
    \label{thm:mpo_arbitrary_periodic_ground_line}
    \lean{MPSTensor.MPOSymmetry.mpv_arbitraryPhysicalMixedInterpolation_ne_zero,
        MPSTensor.MPOSymmetry.arbitraryPhysicalMixedInteraction_groundSpace_eq_span_mpv,
        MPSTensor.MPOSymmetry.arbitraryPhysicalMixedInteraction_groundSpace_finrank}
    \leanok
    \uses{def:mpo_arbitrary_physical_interaction,def:injective,def:mpv}
    Suppose $D_0,D_1>0$ and $A_0,A_1$ are one-site injective.
    For every $\gamma\in[0,1]$ and $N\geq2$, the vector
    $\ket{V^{(N)}(\widetilde A(\gamma))}$ is nonzero, and
    \begin{align}
        \ker\widetilde H_{\gamma,N}
         & =\spn_{\C}\{\ket{V^{(N)}(\widetilde A(\gamma))}\},
         & \dim\ker\widetilde H_{\gamma,N}                    & =1.
        \label{eq:mpo_arbitrary_ground_line}
    \end{align}
\end{theorem}
\begin{proof}\leanok
    \uses{thm:mpo_polar_support_injective,
        thm:mpo_mixed_periodic_line_continuity,
        thm:mpo_periodic_interaction_coordinates,
        thm:mpo_arbitrary_interaction_continuity,
        thm:mpo_isometry_positive_kernel}
    The support-coordinate theorem gives the nonzero periodic vector
    of $\mathcal B(\gamma)$ and its exact ground line.
    The coordinate identity identifies its Euclidean Hamiltonian with
    $H_{\gamma,N}$. Isometric kernel transport gives
    $\ker\widetilde H_{\gamma,N}=J^{\otimes N}(\ker H_{\gamma,N})$.
    Physical mixing takes the support-coordinate vector to
    $\ket{V^{(N)}(\widetilde A(\gamma))}$. Since $J^{\otimes N}$ is
    an isometry, this vector remains nonzero. Its span has dimension one.
\end{proof}

\begin{theorem}[Uniform gap for arbitrary original physical alphabets]
    \label{thm:mpo_arbitrary_uniform_path_gap}
    \lean{MPSTensor.MPOSymmetry.exists_uniform_arbitraryPhysicalMixed_periodic_path_gap}
    \leanok
    \uses{def:mpo_arbitrary_physical_interaction,def:injective}
    If $D_0,D_1>0$ and $A_0,A_1$ are one-site injective, there is a
    constant $0<\varepsilon\leq1$, depending only on these tensors,
    such that for every $\gamma\in[0,1]$, every $N\geq2$, and every
    $v\perp\ker\widetilde H_{\gamma,N}$,
    \begin{align}
        \varepsilon\|v\|\leq\|\widetilde H_{\gamma,N}v\|.
        \label{eq:mpo_arbitrary_uniform_gap}
    \end{align}
\end{theorem}
\begin{proof}\leanok
    \uses{thm:mpo_polar_support_injective,
        thm:mpo_mixed_uniform_periodic_path_gap,
        thm:mpo_periodic_interaction_coordinates,
        thm:mpo_arbitrary_interaction_continuity,
        thm:mpo_isometry_positive_gap}
    Apply the uniform mixed-path gap theorem to the derived injective
    tensors $B_0,B_1$, which have square physical alphabets.
    It gives $\delta>0$, independent of $\gamma\in[0,1]$ and
    $N\geq2$, with $\delta\|x\|\leq\|H_{\gamma,N}x\|$ on the
    kernel complement. The coordinate identity transfers precisely
    this bound to the matrix Hamiltonian. Each $h_\gamma$ is positive,
    so Theorem~\ref{thm:mpo_isometry_positive_gap} applies to the fixed
    inclusion $J$ and gives (\ref{eq:mpo_arbitrary_uniform_gap}) with
    $\varepsilon=\min\{\delta,1\}>0$. Both constants are obtained
    from the endpoint tensors; a support-path gap is not an assumption.
\end{proof}

The resulting continuous positive interactions connect the exact periodic
endpoint vectors in the orthogonal full embeddings
(\ref{eq:mpo_arbitrary_full_embeddings}), with a uniform gap including
$N=2$. The endpoint interactions use the extended limiting supports;
canonical-parent equality has been asserted only in the open interval.
Attaching paths to canonical endpoint interactions, proving covariance
under a common MPO symmetry, and the phase classification with degenerate
ground spaces are separate assertions.
